\scanpage{62}
\mspage{61}

\uline{Notion de support.}

Soit $\unsure{\mathscr{F}|\mathcal{F}}$ un atelier. Soit $\Sigma_{\mathscr{F}}\subset\unsure{\mathcal{P}|P}(\mathscr{F})$ l'ens. des parties de $\mathscr{F}$ qui sont fermées \deleted{à la fois} pour \deleted{$\leq$ et} $\ll$\inserted{ (donc aussi pour $\leq$)}, et stables par Sup pour $\leq$ (Sup finis ou infinis). Soit $\Sigma_M$ l'ens des parties de $M$ fermées pour $\ll$. Alors on a deux bijections inverses l'une de l'autre

\noindent $(\unsure{\alpha|1|*})$
\diagram{\begin{tikzcd}[ampersand replacement=\&,column sep=large,row sep=small]
\Phi \arrow[r,mapsto] \& \Phi\cap M\,\deleted{\unsure{\subset|\in}} \\
\Sigma_{\mathscr{F}} \arrow[r,shift left=1.1ex] \& \Sigma_{\deleted{\unsure{\mathscr{F}}}M} \arrow[l,shift left=1.1ex] \\
\Phi_S=\{F\in\mathscr{F}\mid\widetilde F\subset S\} \& \deleted{\Phi_{\unsure{S|c}}}S \arrow[l,mapsto]
\end{tikzcd}}
\margin{\uline{Iso. d'ens.}\newline\uline{ordonnés}}

moyennant quoi on va les identifier, trouvant un ens. $\Sigma$.

Soit $A\subset\mathscr{F}$ une partie quelconque. On définit son \uline{cosupport} par
\[
\operatorname{cosupp}A=\{G\in\mathscr{F}\mid\forall F\in A,\ F\parallel G\}\subset\mathscr{F}
\]
On voit, grâce à $\unsure{\mathrm{AtD}_0|\mathrm{AtD}_o}$, que
\[
\operatorname{cosupp}A\in\Sigma_{\mathscr{F}}
\]
\deleted{et il \unsure{est}}

% margin note written sideways; slanted upward
\margin{On écrit $F\parallel A$ au lieu de $F\in\operatorname{cosupp}A$.}

On \deleted{\unsure{écrit}} pose
\[
\operatorname{supp}A\inserted{\ (\text{ou }\overline A)}=\operatorname{cosupp}(\operatorname{cosupp}A),
\]
\margin{Ne pas confondre $\overline A$ avec l'adh. de $A$ pour la top. de $\mathscr{F}$ associée soit à $\leq$, soit à $\ll$!}

% margin note written sideways; slanted upward
\margin{On écrit $A\parallel B$ pour les relations équivalentes\newline
$(i)\ A\subset\operatorname{cosupp}B$\newline
$(ii)\ B\subset\operatorname{cosupp}A$\newline
\deleted{\unsure{(iii)}}\unsure{(ii')} $\forall X\in A,\ X\parallel B$\newline
$(iii)\ \forall Y\in B,\ Y\parallel A$\newline
$(iii')\ \forall X\in A,Y\in B,\ X\parallel Y$.\newline
\unsure{Ser.} équivalent aussi (\unsure{cf|ce} dans le cas \unsure{ensembliste}) $\operatorname{supp}A\cap\operatorname{supp}B=\varnothing$\deleted{$\unsure{\mathscr{F}}$} (ça l'implique en tous cas).}

Il est formel (en termes de la seule donnée d'une relation $\parallel$ dans l'ens abstrait $\mathscr{F}$ -- \unsure{ça} marcherait aussi pour une relation quelc. entre deux ens)
